\section{A noncommuting example and further questions}\label{sec:example}

\subsection{The compact triangle}
Consider
\begin{equation}\label{eq:cubic-family}
 \cX=\{xyz-t(x^3+y^3+z^3)=0\}\subset\PP^2\times\Delta.
\end{equation}
Its central fibre is the reduced triangle of coordinate lines. At
each vertex, the derivative of the defining equation with respect
to $t$ is nonzero; elsewhere on the central fibre one spatial
derivative is nonzero. Hence the total space is smooth near the
central fibre, and a smaller disc gives a smooth total space by
properness.

A singular fibre at $t\ne0$ would have
$yz=3tx^2$, $xz=3ty^2$ and $xy=3tz^2$. No coordinate can vanish
at such a projective point. Multiplication then gives $(3t)^3=1$.
Thus the fibres with $0<|t|<1/3$ are smooth plane cubics and are
connected curves of genus one. Restrict the product K\"ahler form
$\omega_{FS}+(i/2)dt\wedge d\bar t$, with projective lines of
area one. Every central component has area one and every pair
meets in one node. The graph matrix is
\begin{equation}\label{eq:triangle-matrix}
 Q=\begin{pmatrix}2&-1&-1\\-1&2&-1\\-1&-1&2\end{pmatrix},
 \qquad\operatorname{spec}Q=\{0,3,3\}.
\end{equation}

Take the trivial rank-two bundle and initial component metrics
$H_1^0,H_2^0,H_3^0$ from \eqref{eq:example-tuple-intro}. Their
determinants are one and they are positive. Direct multiplication gives
\begin{equation}\label{eq:commutator}
 [H_2^0,H_3^0]=\begin{pmatrix}0&2\\-2&0\end{pmatrix}\ne0.
\end{equation}
If a single complex matrix $g$ made all three congruences
$g^*H_i^0g$ diagonal, then $g^*g=D$ would be positive diagonal.
Writing $g=UD^{1/2}$, with $U$ unitary, would make both
$U^*H_2^0U$ and $U^*H_3^0U$ diagonal. This contradicts
\eqref{eq:commutator}.

The constructed initial metrics $\widehat H_t(0)$ are also not
simultaneously diagonalizable by a fixed frame for small $t$.
Choose one fixed point on each component and transport them to
$X_t$. The three metric values tend to the displayed triple. For
any positive triple $(P,Q,R)$ that can be diagonalized by a common
congruence, $P^{-1}Q$ and $P^{-1}R$ commute, since after the
congruence they are simultaneously diagonal under similarity.
Their commutator for the three transported values tends to
\eqref{eq:commutator}, and so is nonzero for small $t$.
Every holomorphic automorphism of the trivial bundle on a compact
connected smooth fibre is a constant matrix, entrywise by the
maximum principle. The obstruction therefore applies to fixed
holomorphic frames on the fibre.

The finite triple remains noncommuting for a short positive
slow-time interval by continuity. On any closed subinterval bounded
away from zero, Theorem~\ref{thm:main} transfers the same
noncommutation test at the three points to the actual flow for small
$t$. The example consequently contains noncommuting actual dynamics,
not only a nonlinear notation for independent scalar equations.
The trace-free Hermitian multiplicity is three. By
\eqref{eq:triangle-matrix}, its six small positive linearized modes
are each asymptotic to $6\pi/L$, and the remaining positive
linearized spectrum has a uniform gap.

The matrices, their determinants and commutator, and the graph
eigenvalues are checked by the accompanying exact symbolic
computation. The displayed algebra supplies the proof; the
computation is a reproducibility check.

\subsection{What the effective description retains}
The limiting small spectral subspace records componentwise parallel endomorphisms.
Its mass is the sum of component areas times the matrix inner
product. A neck measures the disagreement between its endpoint
metrics, with coefficient fixed by its logarithmic conformal
length. For finite amplitudes this disagreement is a Riemannian
logarithm, rather than a difference of matrices in a chosen linear
coordinate. The same geometric data therefore determine both the
linearized operator and the nonlinear finite evolution.

There is a related finite-dimensional structure in the small
two-form spectrum of nodal threefolds
\cite{wuebben-conifold,wuebben-nodal}. There the finite directions
are cohomology classes lost under smoothing, the energy is built
from exceptional-curve periods, and the mass is a harmonic norm.
Those modes are ordinary two-forms. They are not automatically
endomorphism modes of HYM metric flow. In particular a positive
scalar gap can persist in those examples, so the trivial-bundle
linearization need not inherit their small two-form eigenvalues.

The present theorem identifies a nonlinear evolution determined by the
same geometric masses and capacities as the small linearized operator,
in a setting where the componentwise parallel endomorphisms are explicit.
It complements the
fixed-background bundle-deformation results of
\cite{wuebben-local,wuebben-phase}, but does not make their
regularity constants uniform on a degenerating base.

\subsection{Further directions}
The first extension is to nontrivial bundles and nonidentity gluing.
Parallel endomorphisms on different normalization components must
then be compared through the limiting connections and gluing maps.
Their dimensions can differ, and the finite equation need not be
the same unconstrained product flow. Verification of the actual
bundle geometry and of all surviving modes is necessary.

A second question concerns simultaneous degeneration of the base
and approach to a bundle stability wall. Neck capacities and small
extension energies may become comparable. A common reduction would
have to compute their coupling, rather than simply add independently
derived graph operators. Higher-order node-smoothing rates also
require separate geometric analysis: a local equation $uv=t^\nu$
with $\nu>1$ does not have the smooth total space assumed here.

Finally, the finite-slow-time comparison does not determine the
interchange of degeneration and infinite-time limits. Nor does it
give a general HYM theorem on degenerating threefolds or reconstruct
a smooth space from its eigenvalues. These questions require
additional information about the bundle, local geometry and
nonlinear estimates beyond the result proved here.
